\section{Protobuf Parser}\label{sec:proto-parse}

Following the two-step process outlined in Section~\ref{sec:philosophy}, the
Pollux Protobuf parser uses an intermediate representation to handle Protobuf.

A protobuf message on the wire is a sequence of records.

\vspace{1em}
\begin{definition}[Protobuf Record]
  A Protobuf wire record is a tuple with the field number $\id$, the wire type $w$
  and the payload. The wire type will control how long the payload is.
\end{definition}

\vspace{1em}
\begin{definition}[Protobuf Record Wire Type]
  There are four Protobuf wire types represented in a record:
  \begin{itemize}
  \item \wtvar, whose payload is an unsigned number up to $2^{64} - 1$.
  \item \wtil, whose payload is exactly 8 bytes.
  \item \wtis, whose payload is exactly 4 bytes.
  \item $\wtlen{\ell}$, whose payload is exactly $\ell$ bytes.
  \end{itemize}
\end{definition}

As described in Section~\ref{sec:philosophy}, the first step of parsing a
Protobuf value is creating a sequence of records. Once the records are created,
they're sorted, which will group each record by which field is relevant. Then
one linear pass over the list of sorted records can reconstruct the final
message. This is the pass which references the descriptor, converting \wtvar{}
fields into the final integer or boolean, etc.

If there are multiple records for one field, they are combined using the
$\combine{\tau}$ operation. For most types, including all implicit numeric types
and strings, the right-most record is selected. Interestingly, the related
$\byt$ fields aren't mentioned in the specification~\cite{Encoding}, but we've
elected to have then follow the last write wins rule. Elements of repeated
fields can be broken over multiple records, which are merged by the
operator. Messages also follow the merging semantics as per
Section~\ref{sec:proto}.

% Local Variables:
% citar-bibliography: ("../pollux.bib")
% TeX-master: "../pollux.tex"
% TeX-command-default: "Make"
% jinx-local-words: "cardinality protobuf"
% End:
